\documentclass[../main.tex]{subfiles}

\graphicspath{{../image/}{../figure/}}

\begin{document}

\chapter{多空间维中的定常与自相似解}\label{chapter:ch18-多维定常与自相似解}

如本书前面所指出，多空间维中非线性双曲守恒律方程组的一般理论仍是一片未知领域（terra incognita）。尽管如此，二维空间中一批重要问题目前是可处理的，因为它们容许定常或自相似解，此时自变量的数目归结为两个。

本章先介绍二维空间中标量守恒律的 Riemann 问题。在（两个）自相似变量中所得方程仍保持双曲性。出现的波型相当错综复杂，依数据而定可呈现 $32$ 种不同构型，此处只记录其中两种代表情形。

接下来考察平面等熵气流 Euler 方程组的定常或自相似解。自变量数目再次归结为两个；然而代价是所得方程组不再双曲，而呈椭圆--双曲混合型。在跨声速流动中，椭圆区域与双曲区域之间的边界是一条自由边界，须作为解的一部分确定。近期解决的经典问题——无旋流动中楔面的正规激波反射——将作为这一类型的示例较详细地讨论。

\cref{section:ch18-等距浸入}揭示了一个出人意料的联系：把二维 Riemann 流形等距嵌入三维 Euclid 空间这一经典微分几何问题，与 Chaplygin 气体的定常无旋流动之间的联系。

本章最后讨论非线性弹性动力学中的自相似解，用以建模空化现象。

\subfile{Chapter-18/section-01.tex}
\subfile{Chapter-18/section-02.tex}
\subfile{Chapter-18/section-03.tex}
\subfile{Chapter-18/section-04.tex}
\subfile{Chapter-18/section-05.tex}
\subfile{Chapter-18/section-06.tex}
\subfile{Chapter-18/section-07.tex}
\subfile{Chapter-18/section-08.tex}
\subfile{Chapter-18/section-09.tex}

\end{document}
